\SourceLeaf{255}%
unde facile colligitur, aequationes (5), (6) pro casu nostro complete sufficere, dummodo ea, quae indefinita relinquunt, ei conditioni accommodentur, ut \(r\) et \(\varphi\) pro puncto illo initiali atque punctis ab eo infinite parvum distantibus quadrent.

Ceterum quod attinet ad integrationem ipsam aequationum (5), (6), constat, eam reduci posse ad integrationem aequationum differentialium vulgarium, quae tamen plerumque tam intricatae evadunt. ut parum lucri inde redundet. Contra evolutio in series, quae ad usus practicos, quoties de partibus superficiei modicis agitur, abunde sufficiunt, nullis difficultatibus obnoxia est, atque sic formulae allatae fontem uberem aperiunt, ad multa problemata gravissima solvenda. Hoc vero loco exemplum unicum ad methodi indolem monstrandam evolvemus.
\Article{23}
Considerabimus casum eum, ubi omnes lineae, pro quibus \(p\) constans est, sunt lineae brevissimae orthogonaliter secantes lineam, pro qua \(\varphi=0\), et quam tamquam lineam abscissarum contemplari possumus. Sit \(A\) punctum, pro quo \(r=0\), \(D\) punctum indefinitum in linea abscissarum, \(AD=p\), \(B\) punctum indefinitum in linea brevissima ipsi \(AD\) in \(D\) normali, atque \(BD=q\), ita ut \(p\) considerari possit tamquam abscissa, \(q\) tamquam ordinata puncti \(B\); abscissas positivas assumimus in eo ramo lineae abscissarum, cui respondet \(\varphi=0\), dum \(r\) semper tamquam quantitatem positivam spectamus; ordinatas positivas statuimus in plaga ea, ubi \(\varphi\) numeratur inter 0 et \(180^0\).

Per theorema art. 16 habebimus \(\omega=90^0\), \(F=0\). nec non \(G=1\); statuemus insuper \(\BareRoot{E}=n\). Erit itaque \(n\) functio ipsarum \(p,q\), et quidem talis, quae pro \(q=0\) fieri debet \(=1\). Applicatio formulae in art. 18 allatae ad casum nostrum docet, in \textit{quavis} linea brevissima esse debere \(\mathrm{d}\theta=-\frac{\mathrm{d}n}{\mathrm{d}q}\mathord{.}\mathrm{d}p\), denotante \(\theta\) angulum inter elementum huius lineae atque elementum lineae, pro qua \(q\) constans: iam quum linea abscissarum ipsa sit brevissima, atque pro ea ubique \(\theta=0\), patet, pro \(q=0\) ubique fieri debere \(\frac{\mathrm{d}n}{\mathrm{d}q}=0\). Hinc igitur colligimus, si \(n\) in seriem secundum potestates ipsius \(q\) progredientem evolvatur, hanc habere debere formam sequentem
\[
n=1+fqq+gq^3+hq^4+\text{etc.}
\]
ubi \(f,g,h\) etc. erunt functiones ipsius \(p\), et quidem statuemus
